%********************* PAQUETES ADICIONALES *******************

% Para poder escribir acentos tal cual
% \usepackage{isolatin1}
\usepackage{fancyvrb}

% \newcommand{\caja}[1]{
%  \makebox[8mm]{\rule[-1.5ex]{0pt}{6ex}\scriptsize #1}}
% 
% \newcommand{\cajaa}[1]{
%  \makebox[18mm]{\rule[-1.5ex]{0pt}{6ex}\scriptsize #1}}
% 
% \usepackage{shortvrb}
% \MakeShortVerb{\§}
% \DefineShortVerb{\§}

% Para disponer de más símbolos matemáticos
% \usepackage{amsfonts,amssymb,latexsym}
\usepackage{fancyhdr}

% Para que las órdenes como \chapter o \tableofcontents generen
% los títulos en español.
\usepackage[spanish,activeacute]{babel}

% Una extensión para la definición de nuevos teoremas
% \usepackage{theorem}


\usepackage{color} % para importar dibujos coloreados
\usepackage{rotating} % para usar \begin{sideways} que rota tabla 90 grados 
\usepackage[colorinlistoftodos,textwidth=29mm,spanish]{todonotes} % para las notas de TODO:

%\usepackage[sfdefault]{FiraSans} %% option 'sfdefault' activates Fira Sans as the default text font
%\usepackage[T1]{fontenc}
%\renewcommand*\oldstylenums[1]{{\firaoldstyle #1}}

% https://tug.org/FontCatalogue/sansseriffonts.html
\usepackage[sfdefault]{ClearSans}
\usepackage[T1]{fontenc}

%*********************** FORMATO DE PÁGINA **************************

\voffset          -5mm
\cfoot{}

\pagestyle{fancy}
\fancyhead{}
\fancyhead[RE]{{\small  \nouppercase{\leftmark}}}
\fancyhead[LO]{{\small  \nouppercase{\rightmark}}}
\fancyhead[LE,RO]{\thepage}

\setlength{\headheight}{15pt}


\parskip 1.5ex 

%\setlength{\parskip}{2ex} % despues del parrafo, doble linea


%\renewcommand{\baselinestretch}{1,2}


\voffset -0,2cm

\textheight 21.5cm

\textwidth 14.2cm


\oddsidemargin 1cm % estos dos son para que no varíen tanto

\evensidemargin 0,5cm % los margenes entre pags pares e impares


\let\headwidth\textwidth


%\setcounter{tocdepth}{1}


\setlength{\marginparwidth}{27mm}



%==============================================================================
% Cajas numeradas para incluir código
%==============================================================================
%%%%%%%%%%%%%%%%%%%%%%%%%%%%%%%%%%%%%%%%%%%%%%%%%%%%%%%%%%%%%%%%%%%%%%%%%%%%%%%%%%%%%%%%%%%%%
\newlength{\hsbw}

\newcounter{sesioncount}
\setcounter{sesioncount}{0}

%%%%%%%%%%%%%%%%%%% caja no numerada  %%%%%%%%%%%%%%%%%%%%%%%%%

\newenvironment{sesion*}
  {\begin{flushleft}
   \setlength{\hsbw}{\linewidth}
   \addtolength{\hsbw}{-\arrayrulewidth}
   \addtolength{\hsbw}{-\tabcolsep}
   \addtolength{\hsbw}{-\tabcolsep}
   \begin{tabular}{@{}|c@{}|@{}}\hline
   \begin{minipage}[b]{\hsbw}
   \vspace{2mm}
   \begingroup}
  {\endgroup\vspace{-3mm}
   \end{minipage} \\
   \hline
   \end{tabular}
   \end{flushleft}}

%%%%%%%%%%%%%%%%% caja normal numerada  %%%%%%%%%%%%%%%%%%%%%%
\newenvironment{sesion}
  {\begin{flushleft}
   \setlength{\hsbw}{\linewidth}
   \begin{tabular}{@{}|c@{}|@{}}\hline
   \begin{minipage}[b]{\hsbw}
   \vspace*{4mm}
   \em
   \begingroup}
  {\endgroup \vspace{4mm}
   \end{minipage} \\
   \hline
   \end{tabular}
   \end{flushleft}}

%%%%%%%%%%%%%%%%%% caja abierta por debajo %%%%%%%%%%%%%%
\newenvironment{sesion+}
  {\begin{flushleft}
 % \refstepcounter{sesioncount}
   \setlength{\hsbw}{\linewidth}
   \addtolength{\hsbw}{-\arrayrulewidth}
   \addtolength{\hsbw}{-\tabcolsep}
   \addtolength{\hsbw}{-\tabcolsep}
   \begin{tabular}{@{}|c@{}|@{}}\hline
   \begin{minipage}[b]{\hsbw}
   \vspace*{-.5pt}
   \begin{flushright}
   \rule{0.01in}{.15in}\rule{0.3in}{0.01in}\hspace{-0.35in}
   \raisebox{0.04in}{\makebox[0.3in][c]{\footnotesize \thesesioncount}}
                                       % Si pongo \footnotesize\it
                                       % el numerito sale inclinado
   \end{flushright}
   \vspace*{-.47in}
   \begingroup}
  {\endgroup\vspace{-3mm}
   \end{minipage} \\
   %\hline
   \end{tabular}
   \end{flushleft}}

%%%%%%%%%%%% caja abierta por arriba %%%%%%%%%%%%%%%%%%%%%%%%%
\newenvironment{sesion-}
  {\begin{flushleft}
   \setlength{\hsbw}{\linewidth}
   \addtolength{\hsbw}{-\arrayrulewidth}
   \addtolength{\hsbw}{-\tabcolsep}
   \addtolength{\hsbw}{-\tabcolsep}
   \begin{tabular}{@{}|c@{}|@{}}%\hline
   \begin{minipage}[b]{\hsbw}
   \vspace{2mm}
   \begingroup}
  {\endgroup\vspace{-3mm}
   \end{minipage} \\
   \hline
   \end{tabular}
   \end{flushleft}}

\newenvironment{sesionx}
  {\begin{flushleft}
   \setlength{\hsbw}{\linewidth}
   \addtolength{\hsbw}{-\arrayrulewidth}
   \addtolength{\hsbw}{-\tabcolsep}
   \addtolength{\hsbw}{-\tabcolsep}
   \begin{tabular}{@{}|c@{}|@{}}%\hline
   \begin{minipage}[b]{\hsbw}
   \vspace{2mm}
   \begingroup}
  {\endgroup\vspace{-3mm}
   \end{minipage} \\
   %\hline
   \end{tabular}
   \end{flushleft}}

\newcommand{\boxref}[1]
  {\fbox{\footnotesize \phantom{4}\!\!\!\!\!\! \ref{#1}}}
%Si pongo \footnotesize\it
%el numerito sale inclinado

\newcommand{\boxrefdos}[2]
  {\fbox{\footnotesize \phantom{4}\!\!\!\!\!\! \ref{#1}-\ref{#2}}}
%Si pongo \footnotesize\it
%el numerito sale inclinado

\newcommand{\codeindex}[1]{\index{#1@{\tt #1}}}

\newcommand{\labelindex}[1]{\label{#1}\codeindex{#1}}

%%%%%%%%%%%%%%%%%%%%%%%%%%%%%%%%%%%%%%%%%%%%%%%%%%%%%%%%%%%%%%%%%%%%%%%%%%%%%%%
%% § Verbatim extendido                                                      %%
%%%%%%%%%%%%%%%%%%%%%%%%%%%%%%%%%%%%%%%%%%%%%%%%%%%%%%%%%%%%%%%%%%%%%%%%%%%%%%%

\DefineVerbatimEnvironment{CVerbatim}{Verbatim}
{ commandchars=\\\{\},
  codes={\catcode`$=3\catcode`$=3},
  fontsize=\small,
  fontfamily=verdana,
  frame=single}

%=======================================================================

%%%%%%%%%%%%%MACROS%%%%%%%%%%%%%%%%%%
\newcommand{\te}{\vdash}
\newcommand{\lif}{\rightarrow}
\newcommand{\m}{\models}
\renewcommand{\le}{\exists}
\newcommand{\la}{\forall}
\newcommand{\liff}{\leftrightarrow}
\newcommand{\Lor}{\bigvee}
\newcommand{\Land}{\bigwedge}
%=======================================================================



% FIN DE LINEA
\newcommand{\nl}{\newline}



\newcommand{\dip}{\searrow\!\!\!\!\nwarrow}
\newcommand{\dis}{\nearrow\!\!\!\!\swarrow}
\newcommand{\udp}{\nwarrow}
\newcommand{\ddp}{\searrow}
\newcommand{\uds}{\nearrow}
\newcommand{\dds}{\swarrow}

\newcommand{\ua}{\uparrow}
\newcommand{\da}{\downarrow}

\newcommand{\uda}{\uparrow\!\!\downarrow}


\usepackage{cmtt}
\renewcommand{\ttdefault}{cmtt}

%NUEVOS


